\section{Memory Model}
\label{sec:memory-model}

The memory model answers one question: \emph{what may a read return?}
Messages, clocks, lags and scheduling belong to the execution model and exist
only to implement what follows.

It is built in four steps. A partitioned global address space is assumed
outright; a region is a partition of it; a \gls{cell} is a global address plus
a time coordinate; the \gls{store} maps cells to bytes. Everything
characteristic of DSTM is in the third step.

% ===========================================================================
\subsection{The partitioned global address space}
\label{sec:pgas}

DSTM does not construct its address space; it assumes one, of the kind
provided by any PGAS runtime, and adds a coordinate to it. The assumption is
this.

\begin{axiom}[Partitioned global address space]\label{ax:pgas}
There is a set $\Addr$ of \emph{global addresses} such that
\begin{description}[leftmargin=2.9cm,style=nextline,font=\normalfont\bfseries]
  \item[\textnormal{(A1) Partition}] every global address is a pair
    $\gaddr{p}{a}$ of a rank $p \in \Ranks$ and a \emph{local address}
    $a \in \Addrof{p}$, where $\Addrof{p} \defeq \{0, \dots, S_p - 1\}$ for
    some $S_p > 0$:
    \[
      \Addr \defeq \{\, \gaddr{p}{a} : p \in \Ranks,\ a \in \Addrof{p} \,\} .
    \]
  \item[\textnormal{(A2) Uniqueness}] the pair is the address: $\gaddr{p}{a} =
    \gaddr{p'}{a'}$ iff $p = p'$ and $a = a'$. Local address spaces are
    therefore disjoint as partitions of $\Addr$, and $\Addr$ is their disjoint
    union.
  \item[\textnormal{(A3) Affinity}] $\gaddr{p}{a}$ has affinity to $\proc{p}$:
    the bytes reside there, and $\proc{p}$ alone stores them.
  \item[\textnormal{(A4) One-sidedness}] any process may name any global
    address, and access it without the affine process executing application
    code.
\end{description}
\end{axiom}

\begin{definition}[Region]\label{def:region-pgas}
The \emph{region} of rank $p$ is the partition
$\{\, \gaddr{p}{a} : a \in \Addrof{p} \,\}$ of $\Addr$, identified with
$\Addrof{p}$ where no confusion arises. A \gls{pointer} is a global address;
the API encodes it as a \code{(region, offset)} pair in which the region
identifier is the rank, which is (A1) in C.
\end{definition}

(A1)--(A2) are what your intuition for a PGAS already contains: $p$ selects
the partition, $a$ indexes within it, and the pair names one byte-slot
globally. (A3) is why lags are attached to ranks rather than addresses, and
(A4) is why the interface is \code{put}/\code{get} rather than message
passing. Nothing here is specific to DSTM --- UPC++ \cite{upcxx2019} and
Charm++ \cite{charmpp1993} satisfy all four.

% ===========================================================================
\subsection{Versioning: one extra coordinate}
\label{sec:store-is-a-map}

DSTM's entire departure from \cref{ax:pgas} is to index memory by one further
coordinate.

\begin{definition}[Cell]\label{def:cell-pgas}
$\Cells \defeq \Addr \times \Time$. We write $\cell{p}{\eff}{a}$ for
$(\gaddr{p}{a}, \eff)$.
\end{definition}

\begin{definition}[Memory model]\label{def:memory-model}
The state of a DSTM \gls{session} is a \gls{store}
$\Store : \Cells \to \Bytes$ (\cref{def:store}). A read of
$\cell{q}{\eff}{a}$ returns $\Store(\cell{q}{\eff}{a})$.
\end{definition}

``Versioned memory'' suggests a mutable array plus a log of its past states,
so that version $\eff$ is what the array looked like at time $\eff$. That
reading is wrong, and every subsequent confusion follows from it. Time is a
\emph{coordinate of the cell}, on the same footing as the address, not an
annotation on a location that persists across time. Three properties fix the
consequences.

\begin{definition}[Zero default]\label{def:zero-default}
$\Store$ is total. A cell that no \gls{operation} ever writes holds the byte $0$.
\end{definition}

\begin{definition}[Independence across time]\label{def:independence}
For $\eff \ne \eff'$, the cells $\cell{q}{\eff}{a}$ and $\cell{q}{\eff'}{a}$
are distinct and unrelated. Writing one does not affect the other, and the
value of one is never derived from the other. \devmark{dev:non-inheriting}
\end{definition}

\begin{definition}[Single assignment]\label{def:single-assignment}
Each cell is written at most once over the whole execution. A program in which
two distinct writes target the same cell is not a legal DSTM program
(\cref{res:write-disjoint}).
\end{definition}

Together these put DSTM memory in \emph{dynamic single assignment} form
\cite{feautrier1991} over $\Ranks \times \Time \times \Addr$: every cell is
defined at most once and never redefined. This is the discipline SSA imposes
on scalars \cite{cytron1991ssa}, lifted to a distributed setting with
simulation time as the extra index.

\begin{remark}[The consequence programs must live with]
\label{rem:republish}
Because \glspl{version} do not inherit (\cref{def:independence}), state does
not persist across timesteps by itself: a value must be rewritten at $\eff+1$
to remain visible there. There is no read-modify-write across time and no
``current value'' of an address independent of a timestamp. Treating the store
as a per-timestep mailbox --- what both benchmark applications in
\cite{dstm2026} do --- costs nothing; treating it as long-lived state means
re-publishing.
\end{remark}

\begin{remark}[Allocation is not part of the store]
\label{rem:alloc-orthogonal}
$\Store$ is total over $\Addrof{q}$ whether or not an address is allocated
(\cref{def:allocation}). \Gls{allocation} does not decide which \glspl{cell}
\emph{exist}, but which addresses a program may legally \emph{name}. Keeping
the two apart lets the memory model stay a plain total map.
\end{remark}

\begin{definition}[Addressability]\label{def:addressability}
An operation targeting process $\proc{q}$ over the address range
$[a, a + \ell)$ is \emph{addressable} when that whole range lies within a
single extent of $\allocof{q}$:
\[
  \exists\, h .\quad [a,\, a + \ell) \subseteq \extent{q}{h} .
\]
A DSTM program may only issue addressable operations. Targeting an unallocated
address, or straddling two extents, is undefined: the store still answers, but
the byte returned belongs to another structure or to no structure at all
\devmark{dev:darray-bounds}.
\end{definition}

Addressability is a well-formedness condition, not a synchronisation
restriction --- the analogue of not dereferencing a wild pointer, prior to any
question about time. It belongs in the memory model because it concerns
\emph{which memory exists to be named}, not \emph{when it may be touched}.

\begin{remark}[Absence is not observable]
Zero-fill (\cref{def:zero-default}) makes ``nobody wrote here'' and ``somebody
wrote zero'' indistinguishable; a program needing to tell them apart must
encode presence in the value. In exchange reads never fail for want of data,
and a producer may publish only the cells it changes.
\end{remark}

% ===========================================================================
\subsection{Publication: the version life cycle}
\label{sec:publication}

\Cref{def:memory-model} describes the store as though it existed at once. An
implementation must instead decide, per version, the moment after which no
further write can arrive, since only then may it answer a read. That decision
--- promotion from writable to readable --- is the pivot of the design. Status
(\cref{def:status}) moves monotonically through four values.

\begin{figure}[htbp]
\centering
\resizebox{\textwidth}{!}{%
\begin{tikzpicture}[
  node distance=26mm,
  every node/.style={font=\small},
  st/.style={draw,rounded corners=2pt,minimum width=24mm,minimum height=9mm,
             align=center,fill=resbg,draw=dstmblue,thick},
  >={Stealth[round]},
]
  \node[st,fill=white,draw=dstmgrey] (a)             {$\Absent$\\[-1pt]\scriptsize absent};
  \node[st,right=of a]               (w)             {$\Wr$\\[-1pt]\scriptsize writable};
  \node[st,right=of w]               (r)             {$\Rd$\\[-1pt]\scriptsize readable};
  \node[st,right=of r,fill=devbg,draw=dstmred] (x)   {$\Dc$\\[-1pt]\scriptsize discardable};

  \draw[->,thick] (a) -- node[above,font=\scriptsize]{first write}
                          node[below,font=\scriptsize]{\emph{create}} (w);
  \draw[->,thick] (w) -- node[above,font=\scriptsize]{no write can arrive}
                          node[below,font=\scriptsize]{\emph{publish}} (r);
  \draw[->,thick] (r) -- node[above,font=\scriptsize]{no read can arrive}
                          node[below,font=\scriptsize]{\emph{retire}} (x);
  \draw[->,thick,dashed,draw=dstmred] (x.south) .. controls +(0,-9mm) and +(0,-9mm) ..
        node[below,font=\scriptsize]{\emph{evict} (reclaim storage)} (a.south);

  \node[below=13mm of w,font=\scriptsize,align=center,text=dstmgrey]
       {writes accepted\\reads blocked};
  \node[below=13mm of r,font=\scriptsize,align=center,text=dstmgrey]
       {contents frozen\\reads served};
\end{tikzpicture}}
\caption{The version status automaton. Progression is one-way; the dashed
edge is storage reclamation, not a semantic transition --- an evicted version
is indistinguishable from one that was never created, which is sound precisely
because \emph{retire} certifies that nobody will ask for it again.}
\label{fig:status-automaton}
\end{figure}

\begin{definition}[Publication]\label{def:publication}
A version is \emph{published} when its status reaches $\Rd$. From that instant
its content function $\Vbuf{p}{\eff}$ is immutable.
\end{definition}

Both promotions carry the correctness burden, and each is a claim about the
future:

\begin{description}[leftmargin=1.4cm,style=nextline,font=\normalfont\bfseries]
  \item[\emph{publish}] asserts that \emph{no operation that will ever be
    issued can write to this version}. Getting this wrong loses an update: a
    reader observes a version that a later write was still entitled to modify.
  \item[\emph{retire}] asserts that \emph{no operation that will ever be
    issued can read this version}. Getting this wrong either deadlocks a
    reader --- the version it wants is gone and nothing will recreate it ---
    or, worse, lets a subsequent write resurrect the version and serve
    a spuriously zero-filled snapshot.
\end{description}

Neither claim can be checked directly, because a process cannot see the future
and cannot see its neighbours. Both are instead \emph{deduced} from the
\gls{comm-graph} and the \glspl{known-time}, and the deduction is sound only if the
program honours the restrictions of \cref{ch:restrictions}. The exact guards are
given in \cref{sec:l1-rules}; their soundness is
\cref{thm:write-closure,thm:discard-safety}.

\begin{remark}[Eviction is an optimisation, promotion is not]
Reclaiming a $\Dc$ version changes no observable behaviour, so an
implementation with unbounded memory could skip eviction. \emph{Retire} itself
is not optional, since it licenses eviction. This matters for
\cref{thm:space-bound}: the \gls{footprint} is bounded because retirement is
bounded, not because eviction is eager.
\end{remark}

% ===========================================================================
\subsection{What the memory model does not provide}
\label{sec:non-guarantees}

The model is unusually strong for a distributed memory: a read returns
exactly the byte written, with no coherence protocol, fences or visibility
ordering. That strength is bought by the restrictions placed on the program,
and comes with sharp limits.

\begin{description}[leftmargin=0.9cm,style=nextline,font=\normalfont\bfseries]
\item[No atomics, no read-modify-write.]
  There is no compare-and-swap, no fetch-and-add and no way to read a cell and
  write it based on what was read, at the same effective time. The phase
  discipline of \cref{sec:publication} forbids it structurally: by the time a
  version can be read it can no longer be written.

\item[No ordering between distinct cells.]
  Two writes by the same process to different cells are unordered with respect to
  each other in the store, even when one precedes the other in \gls{program-order}.
  Ordering exists only through effective time. A program needing ``$x$ before
  $y$'' must place them at different \glspl{effective-time}.

\item[No conflict detection.]
  Violating single assignment (\cref{def:single-assignment}) is undefined
  behaviour, not an error. There is no abort, no retry and no
  last-writer-wins rule to fall back on \devmark{dev:write-disjoint}.

\item[No partial-failure model.]
  \Glspl{process} do not fail. Messages are not lost, duplicated or reordered in a way
  that matters. This is an MPI-level assumption inherited wholesale.

\item[No dynamic topology.]
  Neither the rank set, the \gls{comm-graph}, the \gls{region} sizes nor a \glslink{dgraph}{dgraph}'s
  edge set may change within a session (\cref{def:session}).
\end{description}

% ===========================================================================
\subsection{Relation to neighbouring memory models}
\label{sec:model-comparison}

DSTM's store differs from familiar designs in one place: what the version
index means, and therefore what makes a version final.

\begin{table}[htbp]
\centering
\small
\begin{tabularx}{\textwidth}{@{}l >{\raggedright\arraybackslash}X
                                 >{\raggedright\arraybackslash}X
                                 >{\raggedright\arraybackslash}X@{}}
\toprule
 & \textbf{DSTM} & \textbf{MVCC} \cite{bernstein1981}
 & \textbf{Time Warp} \cite{jefferson1985} \\
\midrule
Version index
 & simulation time, chosen by the program
 & commit timestamp, chosen by the system
 & send/receive time, with saved state per event \\[2pt]
Version content
 & independent snapshot, zero-filled
 & delta against the previous committed value
 & full checkpoint of process state \\[2pt]
Write conflict
 & forbidden by contract; undefined if violated
 & detected, resolved by abort and retry
 & tolerated, resolved by rollback and antimessages \\[2pt]
Finality of a version
 & deduced from channel lags and neighbour clocks
 & established at commit
 & never; any version may be undone \\[2pt]
Reclamation
 & lag-bounded retirement
 & garbage collection of unreachable snapshots
 & fossil collection below global virtual time \\[2pt]
Read of unwritten cell
 & zero
 & last committed value
 & value in the checkpointed state \\[2pt]
Cost of drift
 & extra live versions (bounded, \cref{thm:space-bound})
 & longer snapshot chains, more aborts
 & rollback and re-execution \\
\bottomrule
\end{tabularx}
\caption{DSTM against multi-version concurrency control and optimistic
simulation. The decisive row is \emph{finality}: MVCC learns a version is
final by observing a commit that has already happened, Time Warp never learns
it at all, and DSTM \emph{predicts} it from static topology plus a lower bound
on neighbour progress.}
\label{tab:model-comparison}
\end{table}

The contrast with a partitioned global \gls{address-space} such as UPC++
\cite{upcxx2019} is of a different kind. PGAS supplies one-sided \code{put} and
\code{get} over remote memory, exactly as DSTM does, but has no version index
at all: a \code{put} overwrites, and the application is responsible for
establishing that no reader is concurrently observing the location. DSTM's
addition is precisely the time coordinate and the runtime-enforced
\gls{publication} discipline that comes with it. Everything else in the interface is
recognisably PGAS.
